\chapter{Cross-Cloud Service Equivalence}
\index{cross-cloud equivalence}\index{service mapping}\index{multi-cloud}%
This appendix is a living reference that maps conceptually equivalent services and
ideas across the six platforms studied in this book series: Amazon Web Services
(AWS), Microsoft Azure, Microsoft Fabric, Google Cloud Platform (GCP), MongoDB, and
Snowflake. The names differ, but the underlying data-engineering concepts---durable
object storage, tiered lifecycle economics, immutability, managed relational engines,
elastic compute, and disciplined data organization---recur everywhere.

\begin{notebox}
This appendix grows with each chapter. Every new chapter contributes the equivalence
tables introduced in its cross-cloud callout, so by the end of the book you will hold
a single consolidated Rosetta Stone for moving fluently between clouds. Where a clean
one-to-one mapping does not exist, the prose notes the closest analog and the caveat.
\end{notebox}

\section{Object and Data-Lake Storage}
\index{object storage!equivalence}\index{storage classes!equivalence}%
The foundational layer covered in Chapter~1. Every platform offers durable, HTTP-addressable
object storage with tiered pricing, lifecycle transitions, versioning, and write-once-read-many
(WORM) controls.

\begin{center}
\footnotesize
\begin{tabular}{@{}p{2.8cm}p{3.2cm}p{3.2cm}p{3.2cm}@{}}
\toprule
\textbf{Capability} & \textbf{AWS} & \textbf{Azure} & \textbf{GCP} \\
\midrule
Object storage & Amazon S3 & ADLS Gen2 / Blob & Cloud Storage \\
Hot tier & S3 Standard & Hot & Standard \\
Infrequent access & S3 Standard-IA & Cool & Nearline \\
Archive (instant) & Glacier Instant Retrieval & Cold & Coldline \\
Deep archive & Glacier Deep Archive & Archive & Archive \\
Lifecycle rules & Lifecycle policies & Lifecycle management & Object Lifecycle Mgmt \\
Versioning & S3 Versioning & Blob versioning & Object Versioning \\
Immutability (WORM) & Object Lock & Immutable blob storage & Bucket Lock \\
Cross-region copy & CRR / SRR & Object replication & Dual/multi-region + Transfer \\
Encryption at rest & SSE-S3 / SSE-KMS & SSE + CMK (Key Vault) & Google-managed / CMEK \\
\bottomrule
\end{tabular}
\end{center}

\noindent The three remaining platforms approach storage differently.
\textbf{Microsoft Fabric} exposes a single logical lake, \emph{OneLake}, and uses
\emph{shortcuts} to virtualize S3, ADLS Gen2, and Google Cloud Storage in place without copying.
\textbf{Snowflake} abstracts storage entirely: table data lives in cloud object storage as
immutable, columnar \emph{micro-partitions} that the service manages on your behalf.
\textbf{MongoDB} persists documents through the \emph{WiredTiger} storage engine, balancing an
internal cache against the operating-system page cache rather than exposing storage tiers.

\section{Managed Relational Databases}
\index{relational databases!equivalence}%
The focus of Chapter~2 for the three hyperscalers. Each provides a fully managed relational
engine, a cloud-native high-performance variant, high availability, read replicas, and a
migration service.

\begin{center}
\footnotesize
\begin{tabular}{@{}p{2.8cm}p{3.2cm}p{3.2cm}p{3.2cm}@{}}
\toprule
\textbf{Capability} & \textbf{AWS} & \textbf{Azure} & \textbf{GCP} \\
\midrule
Managed relational & Amazon RDS & Azure SQL Database & Cloud SQL \\
Cloud-native engine & Amazon Aurora & SQL Hyperscale & AlloyDB \\
Managed PostgreSQL & RDS / Aurora PostgreSQL & Azure DB for PostgreSQL & Cloud SQL / AlloyDB \\
High availability & Multi-AZ & Zone-redundant HA & Regional HA \\
Read scaling & Read Replicas & Read replicas / geo-replicas & Read Replicas \\
Global / geo & Aurora Global Database & Active geo-replication & Cross-region replicas \\
Serverless compute & Aurora Serverless v2 & Azure SQL Serverless & AlloyDB (elastic) \\
Migration service & DMS + SCT & Azure Database Migration Svc & Database Migration Svc \\
\bottomrule
\end{tabular}
\end{center}

\noindent \textbf{Snowflake}, \textbf{Fabric}, and \textbf{MongoDB} are not managed relational
engines, but they solve overlapping problems: Snowflake and Fabric provide SQL analytics over
warehouse/lakehouse storage, while MongoDB provides a distributed document store whose schema
design (Chapter~2) replaces rigid relational normalization with intentional, access-driven modeling.

\section{Elastic Compute and Analytical Warehouses}
\index{elastic compute!equivalence}\index{virtual warehouse!equivalence}%
Snowflake's Chapter~2 topic---separating compute from storage and scaling it elastically---has a
counterpart in every analytics platform.

\begin{center}
\footnotesize
\begin{tabular}{@{}p{2.8cm}p{3.2cm}p{3.2cm}p{3.2cm}@{}}
\toprule
\textbf{Concept} & \textbf{Snowflake} & \textbf{AWS} & \textbf{Azure / GCP} \\
\midrule
Compute unit & Virtual Warehouse & Redshift RA3 / Serverless & Synapse pool / BigQuery slots \\
Elastic concurrency & Multi-cluster warehouse & Concurrency scaling & Autoscale / on-demand \\
Pause when idle & Auto-suspend / resume & Pause / resume & Auto-pause / serverless \\
Billing unit & Credits & RPU / node-hours & DWU / slot-hours \\
Result reuse & Result cache & Result cache & Result / BI cache \\
\bottomrule
\end{tabular}
\end{center}

\noindent \textbf{Microsoft Fabric} meters all workloads against a shared pool of Capacity Units
(CU) on an F-SKU, an approach closer to a single elastic budget than to per-engine warehouses.
\textbf{MongoDB Atlas} scales through cluster tiers with optional auto-scaling of compute and storage.

\section{Data Organization and the Medallion Pattern}
\index{medallion architecture!equivalence}%
Fabric's Chapter~2 topic---the Bronze/Silver/Gold medallion---is a portable design pattern rather
than a product. Every platform implements the same progressive-refinement idea with its own primitives.

\begin{center}
\footnotesize
\begin{tabular}{@{}p{2.6cm}p{3.4cm}p{3.4cm}p{3.0cm}@{}}
\toprule
\textbf{Layer} & \textbf{Fabric} & \textbf{Data lake (AWS/Azure/GCP)} & \textbf{Snowflake} \\
\midrule
Bronze (raw) & Bronze Lakehouse & \texttt{raw/} prefix (S3/ADLS/GCS) & RAW schema / stage \\
Silver (validated) & Silver Lakehouse & \texttt{validated/} prefix & STAGING schema \\
Gold (curated) & Gold Lakehouse & \texttt{curated/} prefix & ANALYTICS schema \\
Storage format & Delta-Parquet & Parquet / Delta / Iceberg & Micro-partitions \\
\bottomrule
\end{tabular}
\end{center}

\noindent \textbf{MongoDB} realizes the same separation with distinct collections or databases
(for example \texttt{raw\_events}, \texttt{validated}, \texttt{curated}) governed by schema-design
patterns rather than a lake layout.

\section{Document Data Modeling}
\index{document modeling!equivalence}%
MongoDB's Chapter~2 schema-design patterns generalize to every document and wide-column store.

\begin{center}
\footnotesize
\begin{tabular}{@{}p{3.0cm}p{3.2cm}p{3.2cm}p{3.0cm}@{}}
\toprule
\textbf{Concept} & \textbf{MongoDB} & \textbf{Amazon DynamoDB} & \textbf{Azure Cosmos DB} \\
\midrule
Primary access unit & Document / collection & Item / table & Item / container \\
Partitioning & Shard key & Partition key & Partition key \\
One-table modeling & Embedding + patterns & Single-table design & Embedding + partitioning \\
Time-series grouping & Bucket Pattern & Composite sort keys & Bucketing by key \\
Optional attributes & Attribute Pattern & Sparse attributes & Sparse properties \\
\bottomrule
\end{tabular}
\end{center}

\noindent Google Cloud's closest analogs are \textbf{Firestore} and \textbf{Bigtable}. In
\textbf{Snowflake} and \textbf{Fabric}, semi-structured documents are stored in \texttt{VARIANT}
(Snowflake) or JSON columns and queried with path expressions rather than modeled as collections.

\section{Globally Distributed and NoSQL Databases}
\index{NoSQL!equivalence}\index{distributed databases!equivalence}%
Chapter~3 branches by platform into high-scale NoSQL (AWS DynamoDB, Azure Cosmos DB), globally
consistent NewSQL (GCP Cloud Spanner), and advanced document modeling (MongoDB). They all answer
the same question: how to scale a database horizontally across partitions and regions.

\begin{center}
\footnotesize
\begin{tabular}{@{}p{2.8cm}p{2.9cm}p{2.9cm}p{2.9cm}p{2.4cm}@{}}
\toprule
\textbf{Concept} & \textbf{DynamoDB} & \textbf{Cosmos DB} & \textbf{Cloud Spanner} & \textbf{MongoDB} \\
\midrule
Model & Key-value / doc & Multi-model & Distributed SQL & Document \\
Partitioning & Partition key & Partition key & Split by key range & Shard key \\
Throughput unit & RCU / WCU & Request Units (RU) & Nodes / processing units & Cluster tier \\
Secondary indexes & GSI / LSI & Automatic indexing & Secondary indexes & Indexes \\
Change stream & DynamoDB Streams & Change Feed & Change streams & Change streams \\
Expiry & TTL & TTL & (manual) & TTL index \\
Multi-region & Global Tables & Multi-region writes & Multi-region (synchronous) & Zones / global clusters \\
Consistency & Eventual / strong & 5 tunable levels & External (TrueTime) & Tunable read/write concern \\
\bottomrule
\end{tabular}
\end{center}

\noindent \textbf{Snowflake} and \textbf{Fabric} are analytical platforms rather than OLTP key-value
stores; they ingest from these operational databases and store semi-structured payloads in
\texttt{VARIANT}/JSON columns. Google's dedicated NoSQL analogs are \textbf{Firestore} (document)
and \textbf{Bigtable} (wide-column).

\section{Managed Apache Spark}
\index{Apache Spark!equivalence}%
Microsoft Fabric's Chapter~3 topic---managed Spark pools and notebooks---has a first-class
equivalent on every cloud.

\begin{center}
\footnotesize
\begin{tabular}{@{}p{2.8cm}p{3.2cm}p{3.2cm}p{3.2cm}@{}}
\toprule
\textbf{Capability} & \textbf{Fabric} & \textbf{AWS} & \textbf{Azure / GCP} \\
\midrule
Managed Spark & Starter / Custom pools & Glue / EMR / EMR Serverless & Synapse Spark / Databricks; Dataproc \\
Notebook UX & Fabric Notebooks & Glue Studio / EMR Studio & Synapse / Databricks; Vertex AI \\
Serverless option & Autoscaling pools & Glue / EMR Serverless & Synapse serverless; Dataproc Serverless \\
Lake format & Delta on OneLake & Delta / Iceberg / Hudi & Delta; Delta / Iceberg \\
\bottomrule
\end{tabular}
\end{center}

\noindent \textbf{Snowflake} offers \textbf{Snowpark} for DataFrame-style processing, and
\textbf{MongoDB} integrates through the Spark Connector.

\section{Query Acceleration and Materialized Views}
\index{materialized views!equivalence}\index{query acceleration!equivalence}%
Snowflake's Chapter~3 acceleration tools---materialized views, the Search Optimization Service, and
transient/temporary tables---map to comparable features in every warehouse.

\begin{center}
\footnotesize
\begin{tabular}{@{}p{2.8cm}p{3.0cm}p{3.0cm}p{3.4cm}@{}}
\toprule
\textbf{Feature} & \textbf{Snowflake} & \textbf{AWS Redshift} & \textbf{Azure Synapse / GCP BigQuery} \\
\midrule
Materialized views & Materialized Views & Materialized views & Materialized views \\
Point-lookup accel. & Search Optimization Svc & Sort keys / zone maps & (indexing) / search indexes \\
Result reuse & Result cache & Result cache & Result-set cache / BI Engine \\
Ephemeral tables & Transient / temporary & Temp tables & Temp tables \\
\bottomrule
\end{tabular}
\end{center}

\noindent \textbf{Fabric} accelerates reads through Direct Lake and its Warehouse engine, while
\textbf{MongoDB} relies on secondary and compound indexes plus the in-memory working set.

\section{In-Memory Caching and Hybrid Ingestion}
\index{caching!equivalence}\index{in-memory cache!equivalence}%
The AWS and Azure Chapter~4 topic---low-latency caching in front of a durable store, fed by hybrid
(edge-to-cloud) ingestion---has a managed Redis-compatible equivalent on every cloud.

\begin{center}
\footnotesize
\begin{tabular}{@{}p{2.8cm}p{3.0cm}p{3.0cm}p{3.4cm}@{}}
\toprule
\textbf{Capability} & \textbf{AWS} & \textbf{Azure} & \textbf{GCP / others} \\
\midrule
Managed cache & ElastiCache (Redis/Memcached) & Cache for Redis & Memorystore (Redis/Memcached) \\
Real-time API layer & AppSync / API Gateway & API Management / SignalR & API Gateway / Apigee \\
Hybrid ingestion & DataSync / Snow family / IoT Core & Data Box / IoT Hub / Arc & Transfer Appliance / IoT Core \\
Streaming intake & Kinesis / MSK & Event Hubs & Pub/Sub \\
\bottomrule
\end{tabular}
\end{center}

\noindent \textbf{Snowflake} and \textbf{Fabric} rely on result and Direct~Lake caches rather than a
standalone Redis tier, while \textbf{MongoDB} serves hot data from its in-memory working set and
Atlas can pair with an external cache.

\section{Wide-Column and Key-Value NoSQL}
\index{wide-column!equivalence}\index{Bigtable!equivalence}%
GCP's Chapter~4 topic---Cloud Bigtable---is a wide-column, high-throughput, low-latency NoSQL store.
Its peers span the managed and open-source ecosystems.

\begin{center}
\footnotesize
\begin{tabular}{@{}p{2.8cm}p{3.0cm}p{3.0cm}p{3.4cm}@{}}
\toprule
\textbf{Characteristic} & \textbf{GCP Bigtable} & \textbf{AWS} & \textbf{Azure / open source} \\
\midrule
Data model & Wide-column & DynamoDB (key-value/doc) & Cosmos DB (multi-model) \\
HBase-compatible & HBase API & Keyspaces (Cassandra) & Managed Instance for Cassandra \\
Scaling unit & Nodes / clusters & On-demand / provisioned & RU/s (Cosmos) \\
Best fit & Time-series, IoT, ad-tech & Serverless key-value & Global multi-model \\
\bottomrule
\end{tabular}
\end{center}

\noindent \textbf{Snowflake} and \textbf{Fabric} address these workloads analytically rather than as
operational key-value stores; \textbf{MongoDB} covers the document end of the same NoSQL spectrum.

\section{Data Protection: Cloning, Time Travel, and Migration}
\index{cloning!equivalence}\index{time travel!equivalence}\index{point-in-time recovery!equivalence}%
Snowflake's Chapter~4 topic---zero-copy cloning, Time Travel, and Fail-safe---plus MongoDB's
relational-to-document migration map to data-protection and migration tooling across the clouds.

\begin{center}
\footnotesize
\begin{tabular}{@{}p{2.8cm}p{3.0cm}p{3.0cm}p{3.4cm}@{}}
\toprule
\textbf{Capability} & \textbf{Snowflake} & \textbf{AWS / Azure} & \textbf{GCP / Fabric} \\
\midrule
Instant clone & Zero-copy clone & Aurora clone / SQL DB copy & BigQuery table clone; OneLake shortcut \\
Point-in-time recovery & Time Travel & PITR (RDS/SQL DB backups) & PITR; Delta time travel \\
Retention safety net & Fail-safe (7 days) & Automated backups & Backups; Delta VACUUM window \\
Schema migration & --- & DMS / SCT & Database Migration Service \\
\bottomrule
\end{tabular}
\end{center}

\noindent \textbf{MongoDB} migrations use \textbf{Relational Migrator} plus the embed-vs-reference
patterns from Chapter~4, and \textbf{Delta Lake} exposes historical versions through
\texttt{VERSION AS OF} time-travel queries.

\section{Cloud Data Warehouses}
\index{data warehouse!equivalence}%
Chapter~5's topic---decoupled compute over open or managed storage with SQL as the primary
interface---is one of the most portable concepts in data engineering.

\begin{center}
\footnotesize
\begin{tabular}{@{}p{2.8cm}p{3.2cm}p{3.2cm}p{3.2cm}p{3.2cm}@{}}
\toprule
\textbf{Capability} & \textbf{Fabric Warehouse} & \textbf{AWS} & \textbf{GCP} & \textbf{Snowflake} \\
\midrule
Warehouse engine & Synapse Data Warehouse & Amazon Redshift & BigQuery & Snowflake virtual warehouses \\
Bulk load command & \texttt{COPY INTO} & \texttt{COPY} & \texttt{LOAD DATA} / bq load & \texttt{COPY INTO <table>} \\
Storage model & Delta-Parquet in OneLake & Managed (RA3) or Spectrum over S3 & Capacitor columnar storage & Micro-partitions on object storage \\
Compute scaling & Capacity CUs (F SKUs) & Nodes / serverless workgroups & Slots / on-demand & Warehouse size (XS--6XL) \\
Cross-db query & Three-part names across items & Redshift Spectrum / federated & Dataset-qualified names & Three-part names, shares \\
Concurrency model & Shared capacity, smoothing & WLM queues / auto-scaling & Slot reservations & Multi-cluster warehouses \\
\bottomrule
\end{tabular}
\end{center}

\noindent \textbf{Databricks} answers with SQL warehouses (serverless, pro, or classic) over
Unity-Catalog-governed Delta tables, and Azure Synapse dedicated SQL pools are the in-family
predecessor. \textbf{MongoDB} has no warehouse concept; its analytics path runs through the
Atlas SQL Interface or export into a columnar engine. The durable lesson: every platform
separates \emph{where bytes live} from \emph{who computes over them}, and bills the two
differently.

\section{SQL over the Open Lake}
\index{lakehouse!SQL equivalence}\index{SQL analytics endpoint!equivalence}%
Chapter~6's read-only SQL surface over open lake storage is the defining lakehouse pattern,
present on every cloud with different degrees of file-hiding.

\begin{center}
\footnotesize
\begin{tabular}{@{}p{2.8cm}p{3.2cm}p{3.2cm}p{3.2cm}p{3.2cm}@{}}
\toprule
\textbf{Capability} & \textbf{Fabric} & \textbf{AWS} & \textbf{GCP} & \textbf{Snowflake} \\
\midrule
SQL over lake files & SQL analytics endpoint (read-only, auto-generated) & Athena over S3; Redshift Spectrum & BigQuery external/BigLake tables & External tables; Iceberg tables \\
Serverless engine & Yes (per-query on capacity) & Athena (per-query/serverless) & Yes (on-demand or slots) & Virtual warehouses \\
Row-level security & RLS on SQL endpoint objects & Lake Formation row/cell filters & BigQuery row access policies & Row access policies \\
Column-level security & GRANT/DENY on columns & Lake Formation column filters & Policy tags & Column masking policies \\
Open table format & Delta-Parquet in OneLake & Iceberg/Hudi/Delta on S3 & Iceberg via BigLake & Iceberg or native \\
\bottomrule
\end{tabular}
\end{center}

\noindent \textbf{Databricks} SQL over Unity Catalog is the closest in-family analog---the same
\emph{open files plus governed SQL layer} idea, with the catalog rather than an auto-generated
endpoint carrying the governance. \textbf{MongoDB} exposes read SQL through the Atlas SQL
Interface rather than over lake files.

\section{Dimensional Modeling and the Semantic Layer}
\index{dimensional modeling!equivalence}\index{semantic model!equivalence}%
Kimball's method (Chapter~7) predates every cloud; what differs per platform is where the
model lives and which language expresses its measures.

\begin{center}
\footnotesize
\begin{tabular}{@{}p{2.8cm}p{3.2cm}p{3.2cm}p{3.2cm}p{3.2cm}@{}}
\toprule
\textbf{Capability} & \textbf{Fabric} & \textbf{AWS} & \textbf{GCP} & \textbf{Snowflake} \\
\midrule
Model home & Power BI semantic model & Redshift tables + QuickSight datasets & BigQuery + Looker semantic layer (LookML) & Snowflake tables + Snowsight/BI tool models \\
Star-schema DDL & Warehouse/lakehouse Delta tables & Redshift (dist/sort keys) & BigQuery (partition/cluster) & Native tables, clustering \\
SCD2 implementation & Delta \texttt{MERGE} & \texttt{MERGE} / staging upserts & \texttt{MERGE} & \texttt{MERGE} / streams+tasks \\
Measure language & DAX & SQL / QuickSight calculated fields & LookML / SQL & SQL / BI tool calcs \\
\bottomrule
\end{tabular}
\end{center}

\noindent \textbf{Databricks} teams model in Delta tables, implement SCD2 with \texttt{MERGE
INTO}, and serve through SQL warehouses or external BI tools. \textbf{MongoDB} modeling is
document-oriented---embedding versus referencing---a different but equally grain-first
discipline.

\section{BI Serving Modes: Import, Live Query, and Direct Lake}
\index{Direct Lake!equivalence}\index{BI serving!equivalence}%
Chapter~8's Direct Lake is Fabric's distinctive contribution: the semantic engine reads the
same open Delta-Parquet files that Spark and SQL use. Other platforms approximate the goal.

\begin{center}
\footnotesize
\begin{tabular}{@{}p{2.8cm}p{3.2cm}p{3.2cm}p{3.2cm}p{3.2cm}@{}}
\toprule
\textbf{Serving pattern} & \textbf{Fabric} & \textbf{AWS} & \textbf{GCP} & \textbf{Snowflake} \\
\midrule
BI reads open lake files & Direct Lake over OneLake Delta & QuickSight over Athena/Spectrum & Looker/Looker Studio over BigLake & External/Iceberg tables to BI tools \\
In-memory BI copy & Power BI Import & QuickSight SPICE & Looker in-memory caches & Cached results / extracts in BI tool \\
Live query to source & DirectQuery (warehouse/endpoint) & QuickSight direct query & Looker live connection & Live connection to Snowflake \\
Cache/update control & Framing (auto/on-demand) & SPICE refresh schedule & Cache policies & Result cache (24h) + BI caching \\
Performance lever & V-Order + column pruning & Sort keys + Spectrum tuning & Partition/cluster pruning & Clustering + result cache \\
\bottomrule
\end{tabular}
\end{center}

\noindent \textbf{Databricks} serves BI through SQL warehouses (with Unity Catalog lineage) and
AI/BI dashboards; partner BI tools connect live or via extracts. \textbf{MongoDB} serves through
the Atlas SQL Interface or Atlas Charts. The shared insight: \emph{copy for speed versus query
for freshness} is the universal BI trade-off.

\section{Visual Orchestration and Bulk Ingestion}
\index{orchestration!equivalence}\index{pipelines!equivalence}%
Chapter~9's visual, activity-based orchestrator descends directly from Azure Data Factory; every
cloud offers the pattern with a rich connector library.

\begin{center}
\footnotesize
\begin{tabular}{@{}p{2.8cm}p{3.2cm}p{3.2cm}p{3.2cm}p{3.2cm}@{}}
\toprule
\textbf{Capability} & \textbf{Fabric} & \textbf{AWS} & \textbf{GCP} & \textbf{Snowflake} \\
\midrule
Visual orchestrator & Data Factory pipelines & Glue Workflows / Step Functions & Cloud Composer (Airflow) / Data Fusion & Tasks / Snowpark orchestration \\
Bulk copy & Copy activity (90+ connectors) & Glue ETL / DMS & Data Fusion / Datastream & \texttt{COPY INTO} / connectors \\
Control flow & ForEach, If, Until, Switch & Step Functions states & Airflow DAG logic & Task DAGs \\
Parameterization & Parameters, variables, expressions & Step Functions JSONPath & Airflow templating (Jinja) & Task parameters \\
Code-first alternative & Notebooks in pipelines & Lambda / Glue scripts & Composer DAGs as code & Snowpark Python \\
\bottomrule
\end{tabular}
\end{center}

\noindent \textbf{Databricks} orchestrates with Workflows (Lakeflow), parameterizing jobs and
notebook tasks; \textbf{MongoDB} uses Atlas Triggers and Functions. The portable skill is the
metadata-driven mindset: config rows enumerate work, and one parameterized pipeline executes it.

\section{Low-Code Transformation}
\index{low-code ETL!equivalence}\index{Dataflows!equivalence}%
Chapter~10's visual query designer (Power Query / M lineage) has an analyst-facing counterpart
in every ecosystem.

\begin{center}
\footnotesize
\begin{tabular}{@{}p{2.8cm}p{3.2cm}p{3.2cm}p{3.2cm}p{3.2cm}@{}}
\toprule
\textbf{Capability} & \textbf{Fabric} & \textbf{AWS} & \textbf{GCP} & \textbf{Snowflake} \\
\midrule
Visual/low-code ETL & Dataflows Gen2 (Power Query) & Glue DataBrew / Glue Studio & Data Fusion / Dataprep heritage & Snowsight worksheets, dbt, partner tools \\
Output destinations & Lakehouse, warehouse, Azure SQL, more & S3, Redshift, databases & BigQuery, Cloud Storage & Snowflake tables \\
Incremental refresh & Yes (per query, date-column driven) & Job bookmarks (Glue) & Incremental pipelines (Data Fusion) & Streams + tasks, dynamic tables \\
Transformation engine & Mashup engine (M) & Spark (Glue Studio) / visual recipes & Data Fusion (Dataproc) & SQL pushdown / Snowpark \\
Data-quality gating & Notebook assertions / pipeline gates & Glue Data Quality & Dataplex data quality & Snowflake data metric functions \\
\bottomrule
\end{tabular}
\end{center}

\noindent \textbf{Databricks} offers lakehouse-native notebooks and declarative Delta Live
Tables rather than a visual designer; \textbf{MongoDB}'s Compass aggregation builder plays the
analyst-facing role. The portable idea: visual tools lower the authoring barrier, but quality
gates and destinations are what make their output production-grade.

\section{Hybrid Connectivity}
\index{hybrid connectivity!equivalence}\index{gateway!equivalence}%
Chapter~11's gateways answer a universal enterprise requirement: reach data behind corporate
networks without exposing it publicly.

\begin{center}
\footnotesize
\begin{tabular}{@{}p{2.8cm}p{3.2cm}p{3.2cm}p{3.2cm}p{3.2cm}@{}}
\toprule
\textbf{Capability} & \textbf{Fabric} & \textbf{AWS} & \textbf{GCP} & \textbf{Snowflake} \\
\midrule
On-prem relay/agent & On-premises data gateway & DMS agents, DataSync agents, Storage Gateway & Transfer appliance / Database Migration Service & External access via connectors \\
VNet/VPC private path & VNet data gateway & VPC endpoints, PrivateLink & Private Service Connect, VPC peering & PrivateLink / Private Service Connect \\
Outbound-only pattern & Gateway initiates outbound 443 & Agents poll outbound & Agents/connectors outbound & Outbound from customer network \\
HA model & Gateway clusters (2+ nodes) & Multi-AZ services & Regional redundancy & Multi-AZ by default \\
\bottomrule
\end{tabular}
\end{center}

\noindent \textbf{Databricks} reaches private sources through VPC peering and PrivateLink;
\textbf{MongoDB Atlas} through private endpoints and VPC peering. The invariant: \emph{the
private side initiates the connection outbound}, so no inbound firewall rule ever exposes the
database.

\section{Orchestration Discipline and Monitoring}
\index{monitoring!equivalence}%
Chapter~12's dependency graphs, retries, and alerting are the most tool-shaped yet conceptually
stable layer of the stack.

\begin{center}
\footnotesize
\begin{tabular}{@{}p{2.8cm}p{3.2cm}p{3.2cm}p{3.2cm}p{3.2cm}@{}}
\toprule
\textbf{Capability} & \textbf{Fabric} & \textbf{AWS} & \textbf{GCP} & \textbf{Snowflake} \\
\midrule
Orchestrator & Pipelines (master/child) & Step Functions / MWAA (Airflow) & Cloud Composer (Airflow) & Task graphs \\
Dependency edges & On Success/Failure/Completion/Skip & States, catch/retry & Airflow trigger rules & \texttt{AFTER} task deps \\
Monitoring & Monitoring Hub + run history & CloudWatch + Step Functions console & Cloud Monitoring / Airflow UI & Task history, resource monitors \\
Alerting & Teams/Outlook activities, Activator & SNS / EventBridge & Alerting policies & Notifications / alerts \\
Lineage hooks & Fabric lineage view & OpenLineage (MWAA) & OpenLineage (Composer) & Access history \\
\bottomrule
\end{tabular}
\end{center}

\noindent \textbf{Databricks} Workflows and \textbf{MongoDB} Atlas Triggers round out the
family. The portable discipline: \emph{an alert without a runbook is noise; a retry without
idempotence is corruption}.

\section{Streaming Ingestion}
\index{streaming ingestion!equivalence}\index{event streaming!equivalence}%
Chapter~13's durable event buffers with fan-out to analytical and operational sinks exist on
every cloud.

\begin{center}
\footnotesize
\begin{tabular}{@{}p{2.8cm}p{3.2cm}p{3.2cm}p{3.2cm}p{3.2cm}@{}}
\toprule
\textbf{Capability} & \textbf{Fabric} & \textbf{AWS} & \textbf{GCP} & \textbf{Snowflake} \\
\midrule
Event ingestion & Eventstreams (Event Hubs lineage) & Kinesis / MSK (Kafka) & Pub/Sub & Snowpipe Streaming \\
Custom producer endpoint & Custom app source & Kinesis PutRecords / Kafka producer & Pub/Sub publisher & Streaming ingest SDK \\
Stream processing & Eventstream operators, KQL & Kinesis Data Analytics / Flink & Dataflow (Beam) & Streams + tasks / dynamic tables \\
Analytical sink & KQL database (Eventhouse) & Redshift / Athena over S3 & BigQuery streaming & Snowflake tables \\
Lake sink & Lakehouse Delta tables & S3 via Firehose & Cloud Storage & Managed staging \\
Delivery semantics & At-least-once & At-least-once (Kinesis) & At-least-once (default) & Exactly-once ingest option \\
\bottomrule
\end{tabular}
\end{center}

\noindent \textbf{Databricks} ingests streams with Structured Streaming, Auto Loader, and Delta
Live Tables; \textbf{MongoDB} produces streams through change streams. The portable disciplines:
unique event IDs at the producer, dedup at the first queryable boundary, and knowing your
source's replay window.

\section{Time-Series and Log Analytics Engines}
\index{time-series!equivalence}\index{KQL!equivalence}%
Chapter~14's Eventhouse is the Azure Data Explorer (Kusto) lineage inside the SaaS workspace;
columnar, append-optimized telemetry engines exist across the ecosystem.

\begin{center}
\footnotesize
\begin{tabular}{@{}p{2.8cm}p{3.2cm}p{3.2cm}p{3.2cm}p{3.2cm}@{}}
\toprule
\textbf{Capability} & \textbf{Fabric} & \textbf{AWS} & \textbf{GCP} & \textbf{Snowflake} \\
\midrule
Telemetry/log engine & Eventhouse (KQL database) & CloudWatch Logs Insights / OpenSearch / Timestream & Cloud Logging / BigQuery streaming analytics & Snowflake tables + time-series functions \\
Query language & KQL & Logs Insights QL / SQL & SQL / Logging query language & SQL \\
Time-series functions & \texttt{make-series}, \texttt{series\_*} family & Timestream SQL functions & BigQuery window/ML functions & Window functions, ASOF joins \\
Incremental state & Materialized views & Scheduled queries / rollups & BigQuery materialized views & Materialized views, dynamic tables \\
Dedup idiom & \texttt{arg\_max} by id & Window dedup & \texttt{ROW\_NUMBER} dedup & \texttt{ROW\_NUMBER} / \texttt{MERGE} \\
Retention control & Retention policies per table & Log group retention / TTL & Log bucket retention / table TTL & Time Travel + data retention \\
\bottomrule
\end{tabular}
\end{center}

\noindent \textbf{Databricks} approaches telemetry with Structured Streaming plus Delta;
\textbf{MongoDB} with native time-series collections. The portable skills: think in
\emph{append-only streams plus incremental state}, bound your lookback windows, and treat late
arrivals as a policy decision, not a surprise.

\section{Data-Driven Alerting}
\index{alerting!equivalence}\index{Data Activator!equivalence}%
Chapter~15's continuous condition evaluation over live data exists across clouds, though
Fabric's no-code Reflex is unusually approachable.

\begin{center}
\footnotesize
\begin{tabular}{@{}p{2.8cm}p{3.2cm}p{3.2cm}p{3.2cm}p{3.2cm}@{}}
\toprule
\textbf{Capability} & \textbf{Fabric} & \textbf{AWS} & \textbf{GCP} & \textbf{Snowflake} \\
\midrule
Data-driven alerts & Data Activator (Reflex) & CloudWatch Alarms + SNS & Cloud Monitoring alerts & Snowflake Alerts \\
Stream-native trigger & Reflex over Eventstream/KQL & Kinesis Analytics + Lambda & Dataflow + Eventarc & Streams + tasks + notifications \\
BI-visual trigger & Reflex from Power BI visuals & QuickSight threshold alerts & Looker alerts/schedules & Via BI tool \\
Actions & Teams, email, Power Automate & SNS targets, Lambda & Pub/Sub, Cloud Functions, webhooks & Email, webhook, stored proc \\
Evaluation latency & Seconds to minutes & Alarm period granularity & Policy-based & Scheduled evaluation \\
\bottomrule
\end{tabular}
\end{center}

\noindent \textbf{Databricks} alerts evaluate SQL on schedules; \textbf{MongoDB} Atlas triggers
fire on change streams. The portable disciplines: \emph{sustained conditions over spikes,
actions with owners, and alert noise treated as a reliability bug}.

\section{Live Dashboards and Reverse ETL}
\index{real-time dashboards!equivalence}\index{reverse ETL!equivalence}%
Chapter~16 closes the loop twice: data to humans (live dashboards) and data back to systems
(reverse ETL).

\begin{center}
\footnotesize
\begin{tabular}{@{}p{2.8cm}p{3.2cm}p{3.2cm}p{3.2cm}p{3.2cm}@{}}
\toprule
\textbf{Capability} & \textbf{Fabric} & \textbf{AWS} & \textbf{GCP} & \textbf{Snowflake} \\
\midrule
Live dashboards & Real-Time Dashboards (KQL) & QuickSight / Grafana on Timestream & Looker / Grafana on BigQuery streaming & Snowsight / connected BI \\
Auto-refresh BI & Power BI automatic page refresh & QuickSight refresh & Looker persistent derived tables & BI tool refresh \\
Reverse ETL target & Dataverse / CRM via pipelines & AppFlow to SaaS & Workflows / AppSheet writes & Snowflake reverse-ETL partners \\
Upsert semantics & Alternate-key upsert to Dataverse & API-level upsert per target & API-level per target & \texttt{MERGE} / external writes \\
Sync cost driver & CUs + target API throttling & Lambda/AppFlow invocations & Function invocations + API quotas & Warehouse credits + API quotas \\
\bottomrule
\end{tabular}
\end{center}

\noindent \textbf{Databricks} serves live BI from SQL warehouses and AI/BI dashboards;
\textbf{MongoDB} has Charts and Atlas-native dashboards. The transferable insight: both loops
are billed---dashboards at refresh cadence, reverse ETL per API call.

\section{Data Science Environments}
\index{data science!equivalence}\index{notebooks!equivalence}%
Chapter~17's notebook-centered data science over lake storage is the dominant pattern across
platforms; the differences are in how tightly tracking, registry, and storage integrate.

\begin{center}
\footnotesize
\begin{tabular}{@{}p{2.8cm}p{3.2cm}p{3.2cm}p{3.2cm}p{3.2cm}@{}}
\toprule
\textbf{Capability} & \textbf{Fabric} & \textbf{AWS} & \textbf{GCP} & \textbf{Snowflake} \\
\midrule
Notebooks & Fabric notebooks (Spark) & SageMaker Studio / EMR notebooks & Vertex AI Workbench / Colab & Snowflake Notebooks \\
Visual prep & Data Wrangler & Data Wrangler (SageMaker) & Vertex / Colab tooling & Snowsight / worksheets \\
Compute for EDA & Starter pools + custom pools & SageMaker instances / EMR & Vertex runtimes / Dataproc & Snowpark warehouses \\
Feature storage & Lakehouse Delta tables & Feature Store (SageMaker) & Vertex AI Feature Store & Feature Store / tables \\
Experiment tracking & Built-in MLflow & SageMaker Experiments / MLflow & Vertex AI Experiments & MLflow / Snowpark ML \\
\bottomrule
\end{tabular}
\end{center}

\noindent \textbf{Databricks} is the closest relative---Fabric's data-science experience
deliberately mirrors its notebook-first, MLflow-tracked shape, including a Feature Store and
AutoML. \textbf{MongoDB} feeds these platforms through connectors rather than hosting ML. The
portable discipline: \emph{explore on samples, engineer features without leakage, store features
as governed tables}.

\section{Experiment Tracking and Model Registries}
\index{MLflow!equivalence}\index{model registry!equivalence}%
Chapter~18's two pieces of ML infrastructure---tracking and registry---converged on MLflow as
the industry's common denominator: native in Fabric and Databricks, deployable everywhere else.

\begin{center}
\footnotesize
\begin{tabular}{@{}p{2.8cm}p{3.2cm}p{3.2cm}p{3.2cm}p{3.2cm}@{}}
\toprule
\textbf{Capability} & \textbf{Fabric} & \textbf{AWS} & \textbf{GCP} & \textbf{Snowflake} \\
\midrule
Experiment tracking & Built-in MLflow & SageMaker Experiments / MLflow on EC2 & Vertex AI Experiments & MLflow (external or integrated) \\
Model registry & Fabric model registry (MLflow) & SageMaker Model Registry & Vertex AI Model Registry & Snowflake Model Registry \\
Autologging & \texttt{mlflow.autolog()} & Same library & Same library & Same library \\
Training compute & Fabric Spark pools & SageMaker Training / EMR & Vertex Training / Dataproc & Snowpark ML / warehouses \\
Artifact store & OneLake (via MLflow) & S3 & Cloud Storage & Stages / tables \\
\bottomrule
\end{tabular}
\end{center}

\noindent \textbf{Databricks}' managed MLflow---with Unity-Catalog-backed registry---is the
direct ancestor of Fabric's integration. \textbf{MongoDB} does not host training; it supplies
features. The portable claim: \emph{if it is not tracked, it did not happen; if it is not
registered, it cannot be governed; if it cannot be rolled back, it is not production}.

\section{AI Enrichment at Scale}
\index{AI enrichment!equivalence}\index{SynapseML!equivalence}%
Chapter~19's managed-AI calls from distributed pipelines have the same shape everywhere: batch
the calls, vault the keys, land the results in the analytical store.

\begin{center}
\footnotesize
\begin{tabular}{@{}p{2.8cm}p{3.2cm}p{3.2cm}p{3.2cm}p{3.2cm}@{}}
\toprule
\textbf{Capability} & \textbf{Fabric} & \textbf{AWS} & \textbf{GCP} & \textbf{Snowflake} \\
\midrule
Distributed ML library & SynapseML on Spark & SageMaker SDK / EMR MLlib & Vertex AI / Dataproc MLlib & Snowpark ML \\
Managed LLM calls & Azure OpenAI via SynapseML & Bedrock from Lambda/EMR & Vertex AI Gemini APIs & Cortex / external functions \\
Cognitive AI services & Azure AI Services & Comprehend, Rekognition, Textract & Vision, Natural Language, Document AI & Cortex AI functions \\
Secret management & Key Vault / workspace identity & Secrets Manager & Secret Manager & External access integrations \\
AI-assisted coding & Copilot in Fabric & CodeWhisperer/Q Developer & Gemini Code Assist & Copilot integrations \\
\bottomrule
\end{tabular}
\end{center}

\noindent \textbf{Databricks} runs the same SynapseML library---it originated in the Synapse
ecosystem---and offers Mosaic AI for LLM and agent work; \textbf{MongoDB} Atlas Vector Search
occupies the adjacent retrieval niche (as does Mosaic AI Vector Search). The portable rules:
\emph{batch instead of row-by-row, vault instead of hardcode, land enrichment in Gold}.

\section{The MLOps Loop}
\index{MLOps!equivalence}\index{batch scoring!equivalence}%
Chapter~20's train--register--score--monitor--retrain loop exists on every platform; what
varies is how much of it the platform holds natively.

\begin{center}
\footnotesize
\begin{tabular}{@{}p{2.8cm}p{3.2cm}p{3.2cm}p{3.2cm}p{3.2cm}@{}}
\toprule
\textbf{Capability} & \textbf{Fabric} & \textbf{AWS} & \textbf{GCP} & \textbf{Snowflake} \\
\midrule
Batch scoring & Spark UDF / \texttt{PREDICT} over Delta & SageMaker Batch Transform & Vertex batch prediction & Snowpark ML / Cortex \\
Orchestrated inference & Pipelines invoking notebooks & Step Functions + SageMaker & Vertex Pipelines & Tasks \\
Drift monitoring & Notebook/pipeline distribution checks & SageMaker Model Monitor & Vertex Model Monitoring & Snowflake model monitors \\
Retrain trigger & Pipeline on drift signal & EventBridge + pipeline & Eventarc + pipeline & Tasks + alerts \\
Model source of truth & MLflow registry (Ch.~18) & SageMaker registry & Vertex registry & Snowflake registry \\
Scoring audit & \texttt{model\_version} on every row & Job metadata & Prediction metadata & Versioned functions \\
\bottomrule
\end{tabular}
\end{center}

\noindent \textbf{Databricks} Mosaic ML and MLflow cover the same loop natively, adding
Lakehouse Monitoring for drift; \textbf{MongoDB} operationalizes predictions as documents
written back by these pipelines. The portable laws: \emph{score with the same features you
trained on, stamp every prediction with its model version, monitor distributions not just
errors}.

\section{Layered Access Control}
\index{access control!equivalence}\index{row-level security!equivalence}%
Chapter~21's three layers---platform roles, resource permissions, in-engine row/column
security---are structurally identical across clouds; only the names change.

\begin{center}
\footnotesize
\begin{tabular}{@{}p{2.8cm}p{3.2cm}p{3.2cm}p{3.2cm}p{3.2cm}@{}}
\toprule
\textbf{Layer} & \textbf{Fabric} & \textbf{AWS} & \textbf{GCP} & \textbf{Snowflake} \\
\midrule
Platform container roles & Workspace roles (Admin/Member/Contributor/Viewer) & IAM policies on resources & IAM roles on projects/folders & Account/database roles \\
Item-level grants & Per-item sharing and permissions & Per-resource policies & Per-resource IAM & Object grants \\
Data-plane scoping & OneLake data access roles & Lake Formation & BigQuery authorized views/Dataplex & Row/column policies \\
Row-level security & SQL RLS predicates & Lake Formation cell filters & Row access policies & Row access policies \\
Identity & Microsoft Entra ID & IAM / Identity Center & Cloud Identity & IdP via SCIM/OAuth \\
Audit & Workspace/tenant logs, Purview & CloudTrail & Cloud Audit Logs & Access history \\
\bottomrule
\end{tabular}
\end{center}

\noindent \textbf{Databricks} governs through Unity Catalog's three-level namespace
(catalog/schema/table) with row filters and column masks; \textbf{MongoDB} Atlas through
database roles and client-side field-level encryption. The portable law: \emph{grant at the
narrowest scope that lets the work happen, and audit the drift between policy and reality}.

\section{Metadata Governance and Erasure}
\index{governance!equivalence}\index{Purview!equivalence}\index{GDPR erasure!equivalence}%
Chapter~22's catalog, classification, lineage, and compliance workflows form a distinct
platform layer in every cloud, increasingly unified across the estate.

\begin{center}
\footnotesize
\begin{tabular}{@{}p{2.8cm}p{3.2cm}p{3.2cm}p{3.2cm}p{3.2cm}@{}}
\toprule
\textbf{Capability} & \textbf{Fabric} & \textbf{AWS} & \textbf{GCP} & \textbf{Snowflake} \\
\midrule
Catalog/governance & Microsoft Purview & DataZone / Glue Catalog + Macie & Dataplex & Horizon Catalog \\
Sensitivity labels & Microsoft Information Protection labels on items & Macie classifications & Sensitive Data Protection / policy tags & Object tagging, masking policies \\
Lineage & Fabric lineage view + Purview & Glue/DataZone lineage & Dataplex lineage & Access history / lineage \\
Endorsement & Promoted / Certified & Business glossary approvals & Data product annotations & Listing certification \\
Erasure mechanics & Crypto-shredding, Delta VACUUM coordination & S3 object deletion + Glacier waits & TTLs, crypto-shredding (CMEK) & Time Travel/Fail-safe windows \\
Legal hold & Purview eDiscovery/holds & Legal holds (S3 Object Lock) & Vault holds & Retention + legal processes \\
\bottomrule
\end{tabular}
\end{center}

\noindent \textbf{Databricks} governs through Unity Catalog (with Purview or Unity lineage
integrations); \textbf{MongoDB} through Atlas auditing and field-level encryption. The hard,
portable lesson: \emph{erasure is a propagation problem across caches, history, and
backups---row deletes alone never satisfy it}.

\section{CI/CD, Capacity, and Disaster Recovery}
\index{CI/CD!equivalence}\index{disaster recovery!equivalence}\index{capacity management!equivalence}%
Chapter~23's infrastructure-as-code discipline, staged promotion, and capacity observability
are universal; Fabric's twist is that workspace items are SaaS objects whose source of truth
must be externalized to Git.

\begin{center}
\footnotesize
\begin{tabular}{@{}p{2.8cm}p{3.2cm}p{3.2cm}p{3.2cm}p{3.2cm}@{}}
\toprule
\textbf{Capability} & \textbf{Fabric} & \textbf{AWS} & \textbf{GCP} & \textbf{Snowflake} \\
\midrule
Source control & Workspace Git integration (ADO/GitHub) & CloudFormation/CDK in Git & Terraform/Deployment Manager in Git & Git integration / Terraform provider \\
Staged promotion & Deployment pipelines (Dev/Test/Prod) & Multi-account pipelines & Multi-project pipelines & Multi-database/account promotion \\
Capacity observability & Capacity Metrics app & CloudWatch & Cloud Monitoring & Account usage views \\
Throttling model & CU smoothing + rejection & Per-service quotas & Quotas + rate limits & Warehouse queuing/credits \\
DR for definitions & Git rehydration & IaC redeploy & IaC redeploy & IaC / replication \\
DR for data & Delta time travel; no auto cross-region & Cross-region replication & Multi-region buckets/datasets & Database replication/failover \\
\bottomrule
\end{tabular}
\end{center}

\noindent \textbf{Databricks} externalizes to Git through Repos and Databricks Asset Bundles
(with Terraform for workspace infrastructure); \textbf{MongoDB} Atlas through Terraform and
API-driven automation. The portable disciplines: \emph{definitions live in Git, data recovery
is a separate design, deploy order is a correctness property}.

\section{Certification Landscape}
\index{certification!equivalence}\index{DP-600!equivalents}%
Chapter~24's study discipline transfers even though the services differ: every platform in the
series has a data-engineering certification.

\begin{center}
\footnotesize
\begin{tabular}{@{}p{2.8cm}p{3.2cm}p{3.2cm}p{3.2cm}p{3.2cm}@{}}
\toprule
\textbf{Element} & \textbf{Fabric} & \textbf{AWS} & \textbf{GCP} & \textbf{Snowflake} \\
\midrule
Data-engineering cert & DP-600 (Fabric Analytics Engineer) & Data Engineer Associate (DEA-C01) & Professional Data Engineer & SnowPro Advanced: Data Engineer \\
Adjacent cert & DP-203 lineage / PL-300 (BI) & Data Analytics specialty (paths vary) & ML Engineer & SnowPro Core first \\
Exam style & Scenario MCQ + case studies & Scenario MCQ & Scenario MCQ & Scenario MCQ \\
Heaviest emphasis & Implement/model/serve analytics & Ingestion, storage, processing & Pipelines, storage, ML touchpoints & Data movement, performance \\
Hands-on value & High---labs mirror exam scenarios & High & High & High \\
\bottomrule
\end{tabular}
\end{center}

\noindent \textbf{Databricks} certifies through the Data Engineer Associate and Professional
exams; \textbf{MongoDB} through Developer and DBA tracks. The portable meta-skill: \emph{study
to the skills outline, measure with mocks, remediate by building---not by rereading}.

\section{The Databricks Column: A Consolidated Map}
\index{Databricks!equivalence map}%
Because Fabric's Spark, Delta, and MLflow lineage runs directly through the Databricks
ecosystem, the two platforms deserve their own consolidated mapping---the densest one-to-one
correspondence in this appendix.

\begin{center}
\footnotesize
\begin{tabular}{@{}p{4.6cm}p{5.2cm}p{5.2cm}@{}}
\toprule
\textbf{Domain} & \textbf{Microsoft Fabric} & \textbf{Databricks} \\
\midrule
Lakehouse storage & OneLake Delta-Parquet tables & Delta Lake on cloud object storage, Unity-Catalog-managed \\
Compute & Fabric Spark starter/custom pools on F-SKU capacity & Classic clusters, serverless compute; billed in DBUs \\
Notebooks & Fabric Notebooks with lakehouse attachment & Databricks Notebooks with cluster/serverless attachment \\
SQL warehousing & Synapse Data Warehouse; SQL analytics endpoint & Databricks SQL warehouses \\
Declarative pipelines & Dataflows Gen2; pipeline activities & Delta Live Tables (Lakeflow declarative pipelines) \\
Orchestration & Data Factory pipelines (master/child) & Databricks Workflows (Lakeflow jobs) \\
Streaming & Eventstreams + Eventhouse (KQL) & Structured Streaming + Auto Loader + Delta Live Tables \\
Experiment tracking & Built-in MLflow (experiments, models) & Managed MLflow with Unity Catalog registry \\
AI services & SynapseML + Azure OpenAI & Mosaic AI; SynapseML also runs on Databricks \\
Governance catalog & Microsoft Purview + workspace/item permissions & Unity Catalog (catalog/schema/table) \\
Source control / CI & Workspace Git integration + deployment pipelines & Repos + Databricks Asset Bundles + Terraform \\
BI serving & Power BI with Direct Lake & AI/BI dashboards and Genie; partner BI via SQL warehouses \\
Certification & DP-600 (Fabric Analytics Engineer) & Data Engineer Associate / Professional \\
\bottomrule
\end{tabular}
\end{center}

\begin{tipbox}
Use this appendix as a study aid for the book's lockstep, cross-cloud approach: when you learn a
concept on one platform, immediately locate its equivalent here and predict how the other six
would express the same idea. That habit turns seven separate vocabularies into one mental model.
\end{tipbox}
